\scanpage{229}
Dans le langage de Deligne des foncteurs dérivés de $\Lambda^i$, compte tenu de $(\Phi_*/k_*)^{\unsure{!|1}}=k[1]$, l'assertion s'écrirait:
\[
 (L\Lambda^i)(k[1])\xrightarrow{\sim}\Lambda^i k[1]=k[+i]
\]
-- plus gén.$^{\mathrm{t}}$ on a, sauf erreur, pour un complexe de \uline{chaînes} de $k$-modules $L_\bullet$ projectif en chaque degré
\[
 (L\Lambda^i)(L_\bullet)\xrightarrow{\sim}\Lambda^i(L_\bullet).
\]
% slanted
\margin{est-ce vrai ?}

On peut aussi noter que la filtration de $\Phi_*$ par
\[
 0\to k_*\to\Phi_*\to\underset{\substack{\rotatebox{90}{$=$}\,\mathrm{défn}\\\Phi^{\unsure{*}}/k_*}}{\Psi_*}\to0
\]
définit pour tt $i\geq0$ une filtration
\[
 0\to\underset{\substack{\rotatebox{90}{$=$}\\\Lambda^{i-1}\Psi_*}}{\Lambda^{i-1}\Psi_*\otimes k_*}
 \to\underset{\text{flasque}}{\Lambda^i\Phi_*}\longrightarrow\Lambda^i\Psi_*\to0
\]
\deleted{d'où} \deleted{\unsure{.}}
\deleted{\[
 0\to(\Lambda^{i-1}\Psi_*)^!\to\underset{\text{flasque}}{(\Lambda^i\Phi_*)^!}
 \longrightarrow\Lambda^i\Psi_*\to0
\]}
d'où la conclusion, car on aura $\pi_j(\Lambda^i\Psi_*)\xleftarrow{\sim}\pi_{j-1}(\Lambda^{i-1}\Psi_*)$ par suite exacte d'homotopie.

\uline{NB} a) Il est probable que $(\Lambda^i\Psi_*)^!$ est \deleted{\unsure{$k\Phi$}} $k[i]$ ($k$ placé en dim $i$). On n'en aura pas besoin.

b) Les \unsure{mêmes} arguments montrent
$\Gamma^i(\deleted{\unsure{\Phi}}\Psi^*)$ \deleted{et}, \deleted{$\unsure{\Lambda}$}$\operatorname{Sym}^i(\Psi^*)$ \inserted{$\otimes^{\unsure{i|*}}\Psi^*$} sont \deleted{des} $K(k,i)$ (mais nettement moins ``économiques'' !)

\uline{Cor. 1.4} Si $A_*$ est un $k$-module $\tfrac12$ simplicial\unsure{?}, on a
\[
 \pi_{\ell}(A_*\otimes\Lambda^i\Psi_*)\simeq\pi_{\ell-i}(A_*)
\]
[de façon plus précise, \inserted{on doit avoir} $A_*\otimes\Lambda^i\Psi_*\simeq\deleted{\unsure{S|\Omega}}\unsure{\Omega|\Sigma|\mathrm{Susp}}^i A_*$ \deleted{\unsure{]}} où $\mathrm{Susp}_a$ est la ``suspension additive'' ]

% archive sheet number: 228

